\documentclass{article}
\usepackage{graphicx}
\usepackage[margin=1.5cm]{geometry}
\usepackage{amsmath}

\begin{document}
\twocolumn

\title{Chapter 1 Warm-Up: Python \texttt{print()} Statements}
\author{Prof. Jordan C. Hanson}

\maketitle

\section{Printing Text in Python}

\begin{enumerate}
\item The Python \texttt{print()} function displays text on the screen.  A string of text is placed inside quotation marks and then inside the parentheses of the \texttt{print()} function.
\begin{itemize}

\item Create a new Python file called \texttt{table\_of\_contents.py}.  Begin by printing the title of the book:
\begin{verbatim}
print("Python Programming: A First Course")
\end{verbatim}
\vspace{0.25cm}

\item Add a second \texttt{print()} statement containing an empty string.  What effect does this have on the output?
\begin{verbatim}
print("")
\end{verbatim}
\vspace{0.25cm}

\item Now print the first entry in the table of contents.  Use periods to connect the chapter title on the left to the page number on the right:
\begin{verbatim}
print("Chapter 1: Getting Started.....1")
\end{verbatim}
\vspace{0.25cm}

\item Run the program.  At this point, your output should contain a book title, a blank line, and the first chapter of the table of contents. \\ \vspace{0.25cm}

\end{itemize}
\end{enumerate}

\section{Warm-Up Exercise}

\begin{enumerate}

\item Complete the program using only simple Python statements and the \texttt{print()} function.  Your program should produce the following table of contents:

\begin{verbatim}
Python Programming: A First Course

Chapter 1: Getting Started....................1
Chapter 2: Variables and Data Types..........17
Chapter 3: Making Decisions..................35
Chapter 4: Repeating with Loops..............57
Chapter 5: Writing Functions.................81
\end{verbatim}

\begin{itemize}

\item Use one \texttt{print()} statement for each line. \\ \vspace{0.25cm}

\item Keep each chapter title on the left and its page number on the right.  Add enough periods between them so that the page numbers align. \\ \vspace{0.25cm}

\item How many \texttt{print()} statements are required to produce the complete table of contents, including the blank line after the book title? Sketch a plan of your code below before beginning in Google CoLab. \\ \vspace{0.5cm}

\item Run your program and compare its output with the example above.  Adjust your code until the output matches the requirement. \\ \vspace{0.5cm}

\end{itemize}
\end{enumerate}

\vspace{4cm}

\section{Name and Student ID}

\begin{enumerate}
\item Please write your name and student ID below, and submit this document to the intructor as you leave class.
\end{enumerate}

\end{document}